% Thesis Table 4.1, generated by scripts/make_master_table.py from
% outputs/data/word_order_summary.tsv. Do not edit by hand.
\providecommand{\podrejeno}[1]{\textcolor{black!70}{#1}}
\begin{table}[htbp]
\centering
\small
\setlength{\tabcolsep}{3.5pt}
\begin{tabular}{@{}lrrrrrrr@{\hspace{1.4em}}r@{}}
\toprule
\textbf{Vir / podkorpus} & $n$ & \textbf{SVO} & \textbf{SOV} & \textbf{VSO} & \textbf{VOS} & \textbf{OSV} & \textbf{OVS} & \shortstack[c]{\textbf{Entro-}\\\textbf{pija}} \\
 & & \% & \% & \% & \% & \% & \% & \\
\midrule
\multicolumn{9}{@{}l}{\textbf{I. Osrednja primerjava}} \\
strokovno (SUK) & 1.105 & \textbf{75,2} & 4,5 & 4,2 & 1,5 & 1,8 & 12,8 & 0,89 \\
publicistično (SUK) & 2.837 & \textbf{71,7} & 6,8 & 3,6 & 1,8 & 1,6 & 14,5 & 0,96 \\
leposlovno (SUK) & 162 & \textbf{69,1} & 4,9 & 3,7 & 2,5 & 2,5 & 17,3 & 1,01 \\
spontani govor (SST) & 168 & \textbf{64,3} & 10,7 & 9,5 & 2,4 & 2,4 & 10,7 & 1,16 \\
\addlinespace[4pt]
\multicolumn{9}{@{}l}{\textbf{II. Drugi analizirani človeški viri}} \\
SSJ (standardno pisno)\textsuperscript{a} & 1.969 & \textbf{71,4} & 6,7 & 4,1 & 1,9 & 1,6 & 14,3 & 0,97 \\
Janes-News & 57.631 & \textbf{72,8} & 5,8 & 5,3 & 3,1 & 2,6 & 10,4 & 0,99 \\
Janes-Blog & 162.867 & \textbf{73,9} & 6,1 & 4,2 & 2,7 & 2,7 & 10,4 & 0,96 \\
Janes-Forum & 159.153 & \textbf{70,2} & 6,0 & 7,7 & 4,6 & 2,8 & 8,8 & 1,07 \\
Janes-Tweet & 496.602 & \textbf{77,5} & 5,8 & 3,5 & 2,3 & 2,1 & 8,9 & 0,86 \\
Janes-Wiki (pogovorne strani) & 11.281 & \textbf{73,0} & 4,0 & 6,7 & 4,4 & 2,5 & 9,4 & 0,99 \\
CLASSLA-Wikipedia (članki) & 296.395 & \textbf{76,0} & 3,6 & 2,5 & 2,2 & 1,5 & 14,2 & 0,84 \\
ParlaMint-SI & 497.572 & \textbf{67,9} & 6,4 & 3,7 & 3,5 & 2,1 & 16,4 & 1,05 \\
KDSP & 48.234 & \textbf{66,6} & 5,1 & 11,0 & 4,4 & 2,5 & 10,4 & 1,13 \\
\addlinespace[4pt]
\multicolumn{9}{@{}l}{\textbf{III. Dopolnilna študija primera: eseji (\textit{Šolar}, \textit{AI-GenT}~1.0)}} \\
človeški eseji & 1.671 & \textbf{74,0} & 6,5 & 7,9 & 0,7 & 2,2 & 8,7 & 0,93 \\
eseji modela GPT-5 & 3.248 & \textbf{89,3} & 3,8 & 1,4 & 0,2 & 0,8 & 4,5 & 0,47 \\
\multicolumn{9}{@{}l}{\hspace{1em}\podrejeno{primerjava modelov (205 esejev na model)\textsuperscript{b}}} \\
\hspace{2em}\podrejeno{GPT-5} & \podrejeno{1.920} & \podrejeno{\textbf{89,0}} & \podrejeno{4,6} & \podrejeno{0,8} & \podrejeno{0,4} & \podrejeno{0,4} & \podrejeno{4,8} & \podrejeno{0,48} \\
\hspace{2em}\podrejeno{gemma-2-27b} & \podrejeno{608} & \podrejeno{\textbf{95,7}} & \podrejeno{0,3} & \podrejeno{1,3} & \podrejeno{0,0} & \podrejeno{0,3} & \podrejeno{2,3} & \podrejeno{0,22} \\
\hspace{2em}\podrejeno{GaMS-27B} & \podrejeno{352} & \podrejeno{\textbf{98,3}} & \podrejeno{0,9} & \podrejeno{0,3} & \podrejeno{0,0} & \podrejeno{0,0} & \podrejeno{0,6} & \podrejeno{0,10} \\
\bottomrule
\end{tabular}
\caption[Porazdelitev besednorednih vzorcev v vseh analiziranih virih]{Porazdelitev šestih besednorednih vzorcev in entropija v vseh analiziranih virih; $n$ označuje število izluščenih parov osebka in predmeta.}
\label{tab:pregled-virov}
\par\smallskip
\begin{minipage}{\linewidth}\footnotesize
Shannonova entropija uporablja naravni logaritem; pri šestih enako pogostih vzorcih doseže največjo vrednost $\ln 6 \approx 1{,}79$. Skladenjske oznake so ročno pregledane v \textit{SSJ} in \textit{SST}, v \textit{SUK} le v delu \mbox{\textit{ssj500k-syn}}; drugod so pripisane samodejno.
\par\smallskip
\textsuperscript{a}~\textit{SSJ} se v veliki meri prekriva z ročno pregledanim delom \textit{SUK} (\mbox{\textit{ssj500k-syn}}), zato rezultati niso neodvisni potrditvi.
\par
\textsuperscript{b}~Po 205 ločeno tvorjenih esejev na podlagi metapodatkov 205 človeških esejev iz nabora 691 parov človek--GPT-5; primerjava modelov zato ni neodvisna ponovitev na drugem gradivu.
\end{minipage}
\end{table}
